\documentclass[a4paper,10pt]{article}
\usepackage[utf8]{inputenc} %pour utiliser tous les caracteres du clavier 
\usepackage[T1]{fontenc} %meme chose ici
\usepackage[francais]{babel} %pour ecrire en francais
\usepackage{listings} %pour citer du code

\usepackage{amsfonts} % pour utiliser les symboles de ensembles (reel...autre)
\usepackage{amsmath} %debut des package pour utiliser les formules de math
\usepackage{amssymb}
\usepackage{mathrsfs}
\usepackage[top=2cm, bottom=2cm, left=2.0cm, right=1.9cm]{geometry}

\usepackage{wrapfig}
\usepackage{graphicx} %pour charger des images
\usepackage[dvipsnames]{xcolor}
\usepackage{float}

\usepackage{textcomp}
\usepackage{url} % pour les url

 \usepackage{listings} % pour ajouter du code
 
\lstset{
language=C++,
basicstyle=\normalsize, % ou ça==> basicstyle=\scriptsize,
upquote=true,
aboveskip={1.5\baselineskip},
columns=fullflexible,
showstringspaces=false,
extendedchars=true,
breaklines=true,
showtabs=false,
showspaces=false,
showstringspaces=false,
identifierstyle=\ttfamily,
keywordstyle=\color[rgb]{0,0,1},
commentstyle=\color[rgb]{0.133,0.545,0.133},
stringstyle=\color[rgb]{0.627,0.126,0.941},
}

\title{NachOS : Multithreading \\ Année 2013-2014}
\author{par Jerôme Barbier, Augustin Husson et David Levayer}
\date{\today}

\begin{document}
  \maketitle

  \begin{center}
    \includegraphics[width=10cm,height=10cm]{robottrading.jpg}\\
    Rapport généré avec \LaTeX
  \end{center}
  \tableofcontents
  \newpage
  \section{Mise en place des threads utilisateurs}
  On va tâcher de permettre à l'utilisateur de créer des threads. Pour cela on met en place les appels systèmes suivants:
  \begin{itemize}
   \item int UserThreadCreate(void f(void *arg), void *arg) : cette fonction permettra de créer un thread et d'exécuter la fonction f dans ce thread
   Retourne -1 en cas d'erreur de création de thread.
   \item void UserThreadExit() : Cette procédure permettra de détruire le thread qu'on a créé précédement
  \end{itemize}
  
  ~~\\
  
  A l'appel système UserThreadCreate, on récupère classiquement les paramètres dans les registres 4 et 5, on appel ensuite la fonction \emph{do\_UserThreadCreate}.
  Cette fonction est disponible dans \emph{userthread.cc}. On va maintenant expliciter ce que contient ce fichier. 
  \subsection{Fichier userthread.cc}
    Dans cette partie ce fichier comporte 3 fonctions: 
    \begin{itemize}
     \item extern int do\_UserThreadCreate(int f, int arg) : cette fonction est appelée par l'appel système \emph{UserThreadCreate}
     \item static void StartUserThread(int f) : cette procédure est appelée par la fonction mère do\_UserThreadCreate
     \item void do\_UserThreadExit() cette procédure est appelée par l'appel système \emph{UserThreadExit}
    \end{itemize}
    
    \subsubsection*{do\_UserThreadCreate}
      Comme on l'a dit, cette fonction est appelée après un appel système. Dans un premier temps, elle va créer un thread :
	\begin{lstlisting}
	 Thread* newThread = new Thread("threadUser");
	\end{lstlisting}
     Afin de pouvoir lui allouer le même espace d'adressage que le processus père, il va falloir faire un fork du programme principal.
     Le problème est que la méthode \emph{Fork} de la classe \emph{Thread}, ne permet pas de passer en paramètre les arguments de la fonction f.
     (On rappel que f est la fonction que souhaite exécuter l'utilisateur dans un thread). C'est pourquoi on met en place la structure suivante:
     \begin{lstlisting}
      struct Serialisation{
	int function; // adresse du pointeur de fonction
	int arg; // adresse du pointeur des arguments
      };
     \end{lstlisting}
    Cette structure est mise dans le fichier \emph{userthread.cc} afin qu'elle reste locale au thread utilisateur.\\

    On crée donc cette structure et on l'initialise avec l'adresse de la fonction \emph{f}, et avec l'adresse des arguments \emph{arg}.
    \begin{lstlisting}
     Serialisation* save = new Serialisation;
     save->function = f;
     save->arg = arg;
    \end{lstlisting}
    On finit par faire l'appel à la méthode \emph{Fork}
    \begin{lstlisting}
     newThread->Fork(StartUserThread,(int)save);
    \end{lstlisting}
    Le fork fait un appel à la procédure StartUserThread. On va donc maintenant parler de cette procédure
    \newpage
    \subsubsection*{StartUserThread}
      Cette procédure permet d'initialiser correctement le thread qu'on a créé dans la fonction précédente. C'est à dire entre autre lui allouer une pile différente de
      celle du processus père.
      
      Dans un premier temps, on s'attache à récupérer notre fonction \emph{f} et ses paramètres \emph{args}. Ce qui se fait avec la ligne ci-dessous:
      \begin{lstlisting}
       Serialisation* restor = (Serialisation*) f;
      \end{lstlisting}
      
      On initialise ensuite les registres du thread, ce qui se fait classiquement en mettant des 0 partouts:
      \begin{lstlisting}
       for(int i=0;i<NumTotalRegs;i++)
	{
		machine->WriteRegister(i,0);
	}
      \end{lstlisting}
      
      On va maintenant positionner le pc sur l'adresse de la fonction \emph{f}. Ceci est fait afin qu'au lancement du thread, la première ligne exécutée 
      soit celle correspondante à \emph{f}.
      \begin{lstlisting}
       machine->WriteRegister(PCReg,restor->function);
      \end{lstlisting}
      Bien évidemment pour que f s'exécute correctement, il est nécessaire de lui fournir ses paramètres qu'elle devra retrouver dans le registre 4.
      \begin{lstlisting}
      machine->WriteRegister(4,restor->arg);
      \end{lstlisting}
      
      la machine nachOS possède un pointeur \emph{pcNext} qui est placé à l'instruction qui suit celle pointée par \emph{pc}. La prochaine instruction 
      se trouve classiquement 4 octets après pc. D'où:
      \begin{lstlisting}
      machine->WriteRegister(NextPCReg,restor->function+4);
      \end{lstlisting}
      
      Il reste maintenant à positionner le pointeur de pile : 
      \begin{lstlisting}
      machine->WriteRegister(StackReg,currentThread->space->BeginPointStack());
      \end{lstlisting}
      Bien évidemment on expliquera par la suite (i.e dans le fichier \emp{addrspace.cc}) comment fonctionne la méthode \emph{BeginPointStack}
      
      Maintenant que tous les registres ont été correctement initialisés, il est temps de lancer le programme!!
      \begin{lstlisting}
       machine->Run();
      \end{lstlisting}
  
 \subsection{Fichier addrspace.cc}
  Afin de positionner le pointeur de pile du nouveau thread, il a fallu modifier le fichier \emph{addrspace.cc et addrspace.h} qui permettent de gérer 
  l'espace mémoir des threads/processus. La principal modification est l'ajout d'un attribut de type \emph{BitMap}
  
  \subsubsection*{L'attribut BitMap}
    Une bitMap est un tableau de bit. Cela permet 
\end{document}
